\documentclass[12pt,a4paper]{article}

% ---------- Pacotes ----------
\usepackage[utf8]{inputenc}
\usepackage[T1]{fontenc}
\usepackage[brazilian]{babel}
\usepackage[margin=2.5cm]{geometry}
\usepackage{amsmath,amssymb}
\usepackage{booktabs,longtable,array}
\usepackage{xcolor}
\usepackage{listings}
\usepackage[hidelinks]{hyperref}

% ---------- Listagens de codigo ----------
\lstset{
  basicstyle=\ttfamily\small,
  breaklines=true,
  frame=single,
  columns=fullflexible,
  extendedchars=true,
  literate={á}{{\'a}}1 {à}{{\`a}}1 {â}{{\^a}}1 {ã}{{\~a}}1
           {é}{{\'e}}1 {ê}{{\^e}}1 {í}{{\'i}}1 {ó}{{\'o}}1 {ô}{{\^o}}1 {õ}{{\~o}}1
           {ú}{{\'u}}1 {ç}{{\c c}}1 {Á}{{\'A}}1 {É}{{\'E}}1 {Ç}{{\c C}}1 {Ã}{{\~A}}1 {Ó}{{\'O}}1
}

% ---------- Notacao COMPARTILHADA (segue o Item 1 corrigido do Fernando) ----------
% Os nomes dos comandos sao capitalizados so para evitar conflito com o LaTeX;
% o texto renderizado sai em minusculas, como no documento do Fernando.
% Se alguem mudar um nome de simbolo, mude AQUI e vale para o documento todo.
\newcommand{\At}{\mathit{at}}
\newcommand{\Lev}{\mathit{lev}}
\newcommand{\Clr}{\mathit{clr}}
\newcommand{\On}{\mathit{on}}
\newcommand{\Cov}{\mathit{cov}}
\newcommand{\Free}{\mathit{free}}
\newcommand{\Stable}{\mathit{stable}}
\newcommand{\Poss}{\mathit{poss}}
\newcommand{\Overlap}{\mathit{overlap}}
\newcommand{\Span}{\mathit{span}}
\newcommand{\Move}{\mathit{move}}
% tamanho do bloco: usar \ell(b)

% ---------- Marcador de pendencia ----------
\newcommand{\todo}[1]{\textcolor{red}{\textbf{[TODO: #1]}}}

\title{Mundo dos Blocos de Tamanho Vari\'avel\\
       \large Representa\c{c}\~ao do Conhecimento e Racioc\'inio Automatizado\\
       Fundamentos de Intelig\^encia Artificial -- Prof. Edjard Mota}
\author{Equipe 5\\
\\Fernando Naum dos Santos Batista, Gabriel Lucas Araújo Matos,
\\Guilherme Paiva Orlando, Luis Felipe Ferreira Cavalcante,
\\Yuri Garcia e Silva}
\date{2 de outubro de 2026}

\begin{document}
\maketitle
\tableofcontents
\newpage

% RESPONSAVEL: gerencia (Guilherme e colega)
\section{Introdu\c{c}\~ao ao Problema}

O Mundo dos Blocos \'e um dom\'inio cl\'assico de planejamento em Intelig\^encia Artificial.
Dado um estado inicial e um estado objetivo, o problema \'e encontrar uma sequ\^encia de a\c{c}\~oes que leve de um ao outro.
Na vers\~ao cl\'assica, todos os blocos t\^em o mesmo tamanho e cada bloco ocupa exatamente uma posi\c{c}\~ao.
Neste trabalho estendemos o dom\'inio para blocos de \emph{comprimentos diferentes}, dispostos sobre uma mesa com pontos discretos e empilh\'aveis em n\'iveis.

\subsection{O dom\'inio}

Temos quatro blocos $a$, $b$, $c$ e $d$, com comprimentos $\ell(a)=\ell(b)=1$, $\ell(c)=2$ e $\ell(d)=3$.
A mesa \'e um segmento com os pontos inteiros de $0$ a $6$, o que d\'a seis \emph{slots} unit\'arios, $s_i=[i,i+1]$ para $i\in\{0,\dots,5\}$.
Um bloco come\c{c}a em um ponto $p$ e cobre $\ell(b)$ slots consecutivos. Os blocos podem ficar sobre a mesa (n\'ivel $0$) ou sobre outros blocos (n\'iveis acima).

Essa extens\~ao traz quatro caracter\'isticas que o caso simples n\~ao tem:
\begin{itemize}
\item \textbf{Tamanhos diferentes:} um bloco pode cobrir mais de um slot e se apoiar em mais de um bloco.
\item \textbf{Espa\c{c}os vazios:} o bloco do n\'ivel de cima pode ficar em ponte sobre dois blocos de baixo, deixando um slot vazio entre eles.
\item \textbf{Movimentos laterais:} al\'em de subir e descer, um bloco pode deslizar na horizontal, ent\~ao a posi\c{c}\~ao em que ele \'e colocado importa.
\item \textbf{Equil\'ibrio:} um bloco maior sobre um menor s\'o \'e aceito se tiver apoio suficiente por baixo. A regra adotada exige apoio em pelo menos $\lceil \ell(b)/2\rceil$ dos slots que o bloco cobre.
\end{itemize}

\subsection{Objetivo do trabalho}

O objetivo \'e propor uma representa\c{c}\~ao do conhecimento para esse dom\'inio e us\'a-la em racioc\'inio automatizado para gerar planos.
O trabalho segue as etapas abaixo, na ordem em que o enunciado pede.
\begin{enumerate}
\item Representar o dom\'inio em L\'ogica de Primeira Ordem (LPO), com os predicados, as leis do dom\'inio e a a\c{c}\~ao $\Move$.
\item Indicar, para cada elemento da representa\c{c}\~ao, o que a a\c{c}\~ao adiciona ou remove (adds e deletes).
\item Usar a representa\c{c}\~ao para obter, manualmente, os planos de cada cen\'ario.
\item Considerar ordem parcial entre as metas, na forma $A \xrightarrow{P} B$.
\item Somente depois, codificar a solu\c{c}\~ao de LPO para l\'ogica proposicional (CNF), gerar os planos com um SAT solver e compar\'a-los com os planos manuais.
\end{enumerate}

\subsection{As tr\^es situa\c{c}\~oes}

O enunciado define tr\^es situa\c{c}\~oes sobre o mesmo conjunto de blocos.
Na Situa\c{c}\~ao 1, parte-se do estado inicial $S_0$ e o plano deve chegar a um dos estados objetivo $S_{f1}$, $S_{f2}$, $S_{f3}$ ou $S_{f4}$.
Na Situa\c{c}\~ao 2, o plano vai de $S_0$ at\'e $S_5$.
Na Situa\c{c}\~ao 3, o plano vai de $S_0$ at\'e $S_7$.
Em todas elas a \'unica a\c{c}\~ao dispon\'ivel \'e $\Move(b,y,p)$, que move o bloco $b$ para cima de $y$ (outro bloco ou a mesa $T$) come\c{c}ando no ponto $p$.

\subsection{Organiza\c{c}\~ao do documento}

A Se\c{c}\~ao 2 apresenta a descri\c{c}\~ao formal do dom\'inio (itens 1, 2 e 4) e a execu\c{c}\~ao manual das tr\^es situa\c{c}\~oes (item 3).
As Se\c{c}\~oes 3 a 7 tratam da codifica\c{c}\~ao em CNF e do uso do SAT solver (item 5): as vari\'aveis e cl\'ausulas, os tr\^es cen\'arios codificados, o mapeamento entre a descri\c{c}\~ao formal e o c\'odigo, a execu\c{c}\~ao passo a passo e a interpreta\c{c}\~ao da sa\'ida do solver.

% RESPONSAVEIS: Fernando (item 1: representacao em LPO e Poss, em 02a)
%               Doze (item 2: adds/deletes, em 02a2)
%               Gerencia (item 4: ordem parcial, em 02c)
\section{Descri\c{c}\~ao Formal do Mundo dos Blocos de Tamanho Vari\'avel}

% RESPONSAVEL: Fernando (item 1) -- convertido do docx Item1_Representacao_LPO_corrigido
\subsection{Representação em Lógica de Primeira Ordem}

\subsubsection{Vocabulário}

\paragraph{Constantes.}
Blocos: $a,b,c,d$. Mesa: $T$. Instantes: $t$.
Pontos $p\in\{0,\dots,6\}$. O slot $s_i$ é o intervalo $[i,i+1]$, com $i\in\{0,\dots,5\}$.
Níveis $l\in\{0,1,2\}$. O nível 0 é a mesa. O nível 2 é o maior que aparece nas três situações.

\paragraph{Função de tamanho.}
\[
\ell(a)=1\quad \ell(b)=1\quad \ell(c)=2\quad \ell(d)=3
\]

\paragraph{Predicados primitivos.}
\begin{center}
\begin{tabular}{p{3.2cm} p{10cm}}
\toprule
Predicado & Significado \\
\midrule
$\At(b,p,t)$  & O bloco $b$ começa no ponto $p$ no instante $t$. Posições válidas: $p\in\{0,\dots,6-\ell(b)\}$. \\
$\Lev(b,l,t)$ & O bloco $b$ está no nível $l$ no instante $t$. \\
\bottomrule
\end{tabular}
\end{center}

\paragraph{Predicados derivados (axiomas de definição).}
\begin{align*}
\Cov(b,s,l,t)       &\leftrightarrow \exists p.\ \At(b,p,t)\land \Lev(b,l,t)\land p\le s<p+\ell(b)\\
\Span(b,p)          &= \{\,s\in\mathbb{N} : p\le s<p+\ell(b)\,\} = [p,p+\ell(b))\\
\Overlap(b,p,y,q)   &\leftrightarrow \Span(b,p)\cap\Span(y,q)\neq\emptyset
                       \quad(\text{isto é, } p<q+\ell(y)\ \land\ q<p+\ell(b))\\
\Free(s,l,b,t)      &\leftrightarrow \neg\exists b'\neq b.\ \Cov(b',s,l,t)\\
\Clr(b,t)           &\leftrightarrow \neg\exists b'\neq b,s,l.\ \Lev(b,l,t)\land\Cov(b,s,l,t)\land\Cov(b',s,l+1,t)\\
\On(b,T,t)          &\leftrightarrow \Lev(b,0,t)\\
\On(b,y,t)          &\leftrightarrow \exists s,l.\ \Cov(b,s,l,t)\land\Cov(y,s,l-1,t)\\
\Stable(b,p,l,t)    &\leftrightarrow l=0\ \lor\ \#\{\,s\in[p,p+\ell(b)) : \exists b'\neq b.\ \Cov(b',s,l-1,t)\,\}\ \ge\ \lceil \ell(b)/2\rceil
\end{align*}

Cada característica pedida no enunciado tem um elemento correspondente:
\begin{center}
\begin{tabular}{p{6.2cm} p{8cm}}
\toprule
Característica & Elemento da LPO \\
\midrule
Tamanhos diferentes & $\ell$ e $\Cov$ (os slots cobertos por cada bloco) \\
Espaços vazios & $\Free$, por slot e por nível \\
Equilíbrio (bloco de cima maior que o de baixo) & $\Stable$ \\
Apoio em dois blocos (ponte) & $\On(b,y_1)\land\On(b,y_2)$, derivado de $\Cov$ \\
Dois blocos lado a lado sobre o mesmo apoio & $\Free$ por slot \\
\bottomrule
\end{tabular}
\end{center}

\subsubsection{Leis do domínio}
Valem em todo estado, para todo instante $t$.
\begin{align*}
&\text{(U1)}\ \forall b,t.\ \exists! p.\ \At(b,p,t) &&\text{unicidade de posição}\\
&\text{(U2)}\ \forall b,t.\ \exists! l.\ \Lev(b,l,t) &&\text{unicidade de nível}\\
&\text{(L)}\ \ \At(b,p,t)\rightarrow 0\le p\le 6-\ell(b) &&\text{limites da mesa}\\
&\text{(E)}\ \ \neg\exists b\neq b',s,l.\ \Cov(b,s,l,t)\land\Cov(b',s,l,t) &&\text{exclusão horizontal}\\
&\text{(S)}\ \ \At(b,p,t)\land\Lev(b,l,t)\rightarrow\Stable(b,p,l,t) &&\text{estabilidade}
\end{align*}

\subsubsection{Ação $\Move(b,y,p)$}
Significado: mover o bloco $b$ para cima de $y$ (bloco ou mesa $T$), começando no ponto $p$.
O nível-alvo é $L=0$ se $y=T$, e $L=l_y+1$ caso contrário, onde $\Lev(y,l_y,t)$. A ação é possível em $t$ quando:
\begin{align*}
\Poss(\Move(b,y,p),t)\leftrightarrow\ &\Clr(b,t)\\
&\land\ y\neq b\ \land\ 0\le p\le 6-\ell(b)\ \land\ L\le 2\\
&\land\ \bigl(y=T\ \lor\ \exists q.\ \At(y,q,t)\land\Overlap(b,p,y,q)\bigr)\\
&\land\ \forall s\in[p,p+\ell(b)).\ \Free(s,L,b,t)\\
&\land\ \Stable(b,p,L,t)\\
&\land\ \neg\bigl(\At(b,p,t)\land\Lev(b,L,t)\bigr)
\end{align*}

Subir, descer e deslizar lateralmente são o mesmo operador: o nível-alvo $L$ e a posição $p$ decidem qual dos três acontece.

\paragraph{Notação.}
$\Move(b,y,p)$ é o termo-ação e o instante $t$ entra apenas em $\Poss(\Move(b,y,p),t)$. No manual e na codificação CNF o instante aparece como quarto argumento, $\Move(b,y,p,t)$ e $\mathit{mv}(b,y,p,t)$; as duas formas são equivalentes. Em $\Span$, $\mathbb{N}$ é o conjunto dos inteiros não negativos (índices de slot).

\paragraph{Desvio em relação ao manual.}
O manual exige $\Clr(y)$, o topo inteiro do bloco de destino livre. Aqui essa exigência foi trocada por $\Free$ por slot, que é o que a pré-condição 5 do manual já pede: nenhum outro bloco ocupa os slots que $b$ vai cobrir no nível-alvo.
Os próprios slides precisam disso. Nas Situações 2 e 3, a transição $S_4\to S_5$ coloca o $b$ sobre o $c$ ao lado do $a$, que já está lá. Com $\Clr(y)$, essas transições são impossíveis. Em $S_{f1}$ e $S_{f2}$ da Situação 1, $a$ e $b$ também ficam lado a lado sobre o mesmo bloco.

\subsubsection{Estados das três situações}
Cada estado é uma conjunção, no instante correspondente, de $\At(b,p)\land\Lev(b,l)$ para os quatro blocos. Escrevemos cada bloco como $b(p,l)$. Por exemplo, $S_0$ da Situação 1 expande para:
\[
\At(c,0,t)\land\Lev(c,0,t)\land\At(a,3,t)\land\Lev(a,0,t)\land\At(b,5,t)\land\Lev(b,0,t)\land\At(d,3,t)\land\Lev(d,1,t)
\]

\paragraph{Situação 1.}
\begin{center}
\begin{tabular}{l l}
\toprule
Estado & Posições $b(p,l)$ \\
\midrule
$S_0$    & $c(0,0)\ a(3,0)\ b(5,0)\ d(3,1)$ \\
$S_{f1}$ & $d(3,0)\ a(4,1)\ b(5,1)\ c(4,2)$ \\
$S_{f2}$ & $d(3,0)\ c(4,1)\ a(4,2)\ b(5,2)$ \\
$S_{f3}$ & $c(0,0)\ a(2,0)\ d(0,1)\ b(5,0)$ \\
$S_{f4}$ & $c(0,0)\ a(0,1)\ d(2,0)\ b(5,0)$ \\
\bottomrule
\end{tabular}
\end{center}

\paragraph{Situação 2.}
\begin{center}
\begin{tabular}{l l}
\toprule
Estado & Posições $b(p,l)$ \\
\midrule
$S_0$ & $c(0,0)\ a(0,1)\ b(1,1)\ d(3,0)$ \\
$S_1$ & $c(0,0)\ a(0,1)\ b(2,0)\ d(3,0)$ \\
$S_2$ & $c(0,0)\ b(2,0)\ a(2,1)\ d(3,0)$ \\
$S_3$ & $b(2,0)\ a(2,1)\ d(3,0)\ c(4,1)$ \\
$S_4$ & $b(2,0)\ d(3,0)\ c(4,1)\ a(4,2)$ \\
$S_5$ & $d(3,0)\ c(4,1)\ a(4,2)\ b(5,2)$ \\
\bottomrule
\end{tabular}
\end{center}

\paragraph{Situação 3.}
\begin{center}
\begin{tabular}{l l}
\toprule
Estado & Posições $b(p,l)$ \\
\midrule
$S_0$ & $c(0,0)\ a(3,0)\ b(5,0)\ d(3,1)$ \\
$S_1$ & $c(0,0)\ a(3,0)\ b(5,0)\ d(0,1)$ \\
$S_2$ & $c(0,0)\ d(0,1)\ b(5,0)\ a(5,1)$ \\
$S_3$ & $c(0,0)\ d(2,0)\ b(5,0)\ a(5,1)$ \\
$S_4$ & $c(0,0)\ d(2,0)\ b(5,0)\ a(0,1)$ \\
$S_5$ & $c(0,0)\ d(2,0)\ a(0,1)\ b(1,1)$ \\
$S_6$ & igual a $S_5$ \\
$S_7$ & $c(0,0)\ d(3,0)\ a(0,1)\ b(1,1)$ \\
\bottomrule
\end{tabular}
\end{center}

\paragraph{Verificação e observações.}
\begin{itemize}
\item Todos os estados acima satisfazem as leis (U1), (U2), (L), (E) e (S). Em $S_0$ da Situação 1, por exemplo, o $d$ cobre os slots 3, 4 e 5 no nível 1 e tem apoio nos slots 3 ($a$) e 5 ($b$): $2\ge\lceil 3/2\rceil$. Daí vem $\On(d,a)\land\On(d,b)$, com o slot 4 vazio embaixo (ponte).
\item No slide da Situação 3, $S_5$ e $S_6$ estão desenhados iguais, e isso é intencional no enunciado. $S_6$ é o mesmo estado de $S_5$ (repetição, sem ação entre eles), e o plano segue de $S_5$ direto para $S_7$.
\item O $S_0$ da Situação 3 é o mesmo $S_0$ da Situação 1.
\end{itemize}


% RESPONSAVEL: Doze (item 2) -- convertido do PDF "Parte 2.1 - Acoes" (versao corrigida).
% Correcoes aplicadas apos a revisao: digitacao e o tratamento dos casos p0=p e l0=L no DELETE.
\subsection{Efeitos das ações e tabela de adds/deletes}

\subsubsection{Ações associadas aos elementos da representação}
Os predicados $\At$ e $\Lev$ são os únicos elementos que são alterados pelo $\Move$. Já os demais elementos são estáticos, como o $\ell(b)$, ou são derivados: $\Cov$, $\Free$, $\Clr$, $\On$ e $\Stable$. Portanto seus valores são obtidos a partir do $\At$, $\Lev$ e $\ell(b)$.

\subsubsection{Relação entre os símbolos e a ação move}
\begin{center}
\begin{tabular}{p{3.2cm} p{3cm} p{3cm} p{5cm}}
\toprule
Símbolo & ADD & DELETE & Uso \\
\midrule
$\ell(b)$ & -- & -- & Propriedade estática usada para limites, cobertura e estabilidade. \\
$\At(b,p,t)$ & Nova posição de $b$ & Posição anterior de $b$ & Posição do bloco. \\
$\Lev(b,l,t)$ & Novo nível de $b$ & Nível anterior de $b$ & Nível do bloco. \\
$\Cov(b,s,l,t)$ & -- & -- & Derivado de $\At$, $\Lev$ e $\ell(b)$. \\
$\Free(s,l,b,t)$ & -- & -- & Pré-condição do destino. \\
$\Clr(b,t)$ & -- & -- & Pré-condição para mover $b$. \\
$\On(b,y,t)$ & -- & -- & Derivado de $\Cov$. \\
$\Stable(b,p,l,t)$ & -- & -- & Testado para garantir a estabilidade do destino. \\
$\Span(b,p)$ & -- & -- & Derivado de $p$ e $\ell(b)$, representando os espaços ocupados pelo bloco. \\
$\Overlap(b,p,y,q)$ & -- & -- & Testado em $\Poss(\Move(b,y,p),t)$ para verificar sobreposição com o apoio. \\
\bottomrule
\end{tabular}
\end{center}

A ação é $\Move(b,y,p)$, onde $b$ é o bloco a ser movimentado, $y$ o destino (bloco ou mesa) e $p$ é a nova posição de $b$.
As condições de possibilidade dessa ação já são definidas por $\Poss(\Move(b,y,p),t)$, incluindo $\Clr$, $\Free$, limites de posições e $\Stable$.

\subsubsection{Efeitos ADD e DELETE da ação}
Se a ação $\Move(b,y,p)$ é executada no instante $t$, o bloco $b$ passa a ocupar uma nova posição e novo nível no instante $t+1$.

\paragraph{ADD.} A nova posição e o novo nível são adicionados:
\[
\At(b,p,t+1)\qquad \Lev(b,L,t+1)
\]
Onde $L=0$ se o destino for a mesa $T$ e, se $y$ for outro bloco com $\Lev(y,l_y,t)$, então $L=l_y+1$.

\paragraph{DELETE.} A posição anterior do bloco $b$ deixa de valer, desde que seja diferente da nova. Se antes da ação:
\[
\At(b,p_0,t)
\]
e $p_0\neq p$, então após o movimento:
\[
\neg\At(b,p_0,t+1)
\]
Se $p_0=p$ (o bloco apenas sobe ou desce no mesmo ponto), o $\At$ não é apagado: $\At(b,p,t+1)$ já vale pelo ADD, e apagá-lo o contradiria.

Da mesma forma, se o nível anterior for:
\[
\Lev(b,l_0,t)
\]
e $l_0\neq L$, então:
\[
\neg\Lev(b,l_0,t+1)
\]
Se $l_0=L$ (o bloco desliza lateralmente no mesmo nível), o $\Lev$ não é apagado.

Portanto, a ação $\Move$ remove diretamente a posição e o nível antigos, quando mudam, e adiciona a nova posição e o novo nível.

\subsubsection{Predicados derivados}
Os predicados derivados são: $\Cov$, $\Free$, $\Clr$, $\On$, $\Stable$.
Após $\At$ e $\Lev$ serem alterados, esses predicados são recalculados pelas definições apresentadas no item 1.
Por exemplo, $\On$ é derivado a partir da posição, nível e sobreposição dos blocos, enquanto $\Clr$ depende da existência ou não de blocos ocupando posições acima.
Isso significa que não é necessário escrever: $\neg\Clr(y,t+1)$.

\subsubsection{Estabilidade após o movimento}
A estabilidade do bloco movimentado deve ser verificada antes que o movimento seja permitido: $\Stable(b,p,l,t)$.
Assim, o novo estado só pode ser obtido se o bloco possuir suporte suficiente na posição e nível de destino.

A estabilidade dos demais blocos não precisa ser verificada novamente porque $\Clr(b,t)$ é uma pré-condição para mover $b$. Pela definição de $\Clr$, isso garante que não existe nenhum bloco ocupando o nível imediatamente acima de $b$ e apoiado sobre sua região. Portanto, ao retirar $b$ de sua posição atual, nenhum outro bloco perde suporte por causa dessa remoção. Assim, basta verificar a estabilidade do próprio bloco movimentado em sua nova posição.

\subsubsection{Aplicações}
As regras acima são gerais e se aplicam às Situações 1, 2 e 3.
Em todos os casos, permanecem os mesmos:
\begin{itemize}
\item os efeitos ADD sobre $\At$ e $\Lev$;
\item os efeitos DELETE da posição e do nível anteriores do bloco movimentado;
\item as pré-condições da ação $\Move$;
\item os predicados derivados $\Cov$, $\Free$, $\Clr$, $\On$ e $\Stable$;
\item os axiomas de persistência dos blocos não movimentados.
\end{itemize}

\paragraph{Persistência.} Para qualquer bloco $b'\neq b$, se a ação executada em $t$ move apenas o bloco $b$, então $b'$ mantém sua posição e seu nível em $t+1$:
\begin{align*}
b'\neq b &\rightarrow \bigl(\At(b',p',t)\leftrightarrow\At(b',p',t+1)\bigr)\\
b'\neq b &\rightarrow \bigl(\Lev(b',l',t)\leftrightarrow\Lev(b',l',t+1)\bigr)
\end{align*}
Assim, a ação $\Move(b,y,p)$ altera diretamente apenas o bloco $b$. Os demais permanecem inalterados.

Apenas o que muda entre as situações são:
\begin{itemize}
\item os estados iniciais;
\item os estados objetivos;
\item as instâncias da ação $\Move$.
\end{itemize}


% RESPONSAVEL: gerencia (Guilherme) -- item 4 (ordem parcial). Implementacao: ordem_parcial.py e grupo (ORD) de gerar_cnf.py.
\subsection{Ordem Parcial}

Um plano linear $A_1,A_2,\dots,A_k$ fixa uma ordem total entre as ações. Um \emph{plano de ordem parcial} fixa apenas as ordens que são necessárias, e deixa livre a ordem das ações independentes. Esta seção formaliza isso para o mundo dos blocos de tamanho variável (itens 1 e 2) e depois trata a ordem parcial entre as \emph{metas}, que é o que entra na codificação em CNF.

\subsubsection{Plano de ordem parcial}

Um plano de ordem parcial é uma tripla $\langle \mathcal{A},\prec,\mathcal{L}\rangle$:
\begin{itemize}
\item $\mathcal{A}$ é um conjunto de passos, cada um uma instância $\Move(b,y,p)$;
\item $\prec$ é uma relação de ordem entre passos: irreflexiva e transitiva, logo sem ciclos. $A_i\prec A_j$ significa que $A_i$ é executado antes de $A_j$;
\item $\mathcal{L}$ é um conjunto de \emph{elos causais} $A_i\xrightarrow{\varphi}A_j$.
\end{itemize}

\paragraph{Fatos.} O estado é dado por $\At$ e $\Lev$, e as pré-condições de $\Move$ são derivadas. Os fatos que ligam uma ação à outra são sobre slots:
\[
\mathit{livre}(s,l)\equiv\Free(s,l,\cdot,t)=\neg\exists b'.\ \Cov(b',s,l,t)\qquad
\mathit{ocupado}(s,l)\equiv\neg\mathit{livre}(s,l)
\]

\paragraph{Pré-condições e efeitos em termos de fatos.} Para $\Move(b,y,p)$, com $b$ antes em $(p_0,l_0)$ e nível-alvo $L$:
\begin{center}
\begin{tabular}{l p{9.2cm}}
\toprule
Pré-condição & Fatos exigidos \\
\midrule
(P1) $\Clr(b,t)$ & $\mathit{livre}(s,l_0+1)$ para todo $s\in\Span(b,p_0)$, porque $\Clr(b,t)\leftrightarrow\forall s\in\Span(b,p_0).\ \Free(s,l_0+1,b,t)$ \\
(P4) destino livre & $\mathit{livre}(s,L)$ para todo $s\in\Span(b,p)$ (exceto os slots que o próprio $b$ já ocupa nesse nível) \\
(P5) estabilidade & $\mathit{ocupado}(s,L-1)$ para os slots de $\Span(b,p)$ que servem de apoio \\
\midrule
Efeito & Produz $\mathit{livre}(s,l_0)$ para $s\in\Span(b,p_0)$ e $\mathit{ocupado}(s,L)$ para $s\in\Span(b,p)$; desfaz os fatos contrários \\
\bottomrule
\end{tabular}
\end{center}

\paragraph{Elo causal.} $A_i\xrightarrow{\varphi}A_j$ significa que $A_i$ produz o fato $\varphi$, $A_j$ exige $\varphi$ como pré-condição e $\varphi$ vale em todo instante entre os dois. O elo impõe $A_i\prec A_j$.

\paragraph{Ameaça e resolução.} Um passo $A_k$ \emph{ameaça} o elo $A_i\xrightarrow{\varphi}A_j$ se desfaz $\varphi$ e pode ser executado entre $A_i$ e $A_j$. Por exemplo, se $A_i$ libera um slot para $A_j$, qualquer ação que ocupe esse slot é uma ameaça. A ameaça é resolvida por \emph{promoção} ($A_k\prec A_i$) ou \emph{rebaixamento} ($A_j\prec A_k$). Em um plano totalmente ordenado não há ameaças, porque a ordem já resolve todas.

\paragraph{Linearização e correção.} Uma linearização é uma ordem total compatível com $\prec$. O plano de ordem parcial é correto se toda linearização é executável (cada ação satisfaz P1 a P6 no seu momento) e termina no estado meta. A execução manual abaixo testa isso por simulação.

\subsubsection{Execução manual: planos do item 3 como ordem parcial}

Para cada plano manual do item 3, calculamos os elos causais pelos efeitos de cada ação (colunas Add e Delete do item 3) e testamos todas as ordens de execução das mesmas ações. O script \texttt{ordem\_parcial.py} automatiza as duas etapas.

\begin{center}
\begin{tabular}{l p{3.6cm} p{8.4cm}}
\toprule
Plano & Passos & Elos causais $A_i\xrightarrow{\varphi}A_j$ \\
\midrule
Sit.\ 1 ($S_{f3}$) & $A_1=\Move(d,c,0)$, $A_2=\Move(a,T,2)$ & $A_1\xrightarrow{\mathit{livre}(3,1)}A_2$ \\
\midrule
Sit.\ 1 ($S_{f4}$) & $A_1=\Move(d,c,0)$, $A_2=\Move(a,b,5)$, $A_3=\Move(d,T,2)$, $A_4=\Move(a,c,0)$ &
$A_1\xrightarrow{\mathit{livre}(3,1),\,\mathit{livre}(5,1)}A_2$,\ $A_2\xrightarrow{\mathit{livre}(3,0)}A_3$,\ $A_3\xrightarrow{\mathit{livre}(0,1)}A_4$ \\
\midrule
Sit.\ 2 ($S_5$) & $A_1=\Move(b,T,2)$, $A_2=\Move(a,b,2)$, $A_3=\Move(c,d,4)$, $A_4=\Move(a,c,4)$, $A_5=\Move(b,c,5)$ &
$A_1\xrightarrow{\mathit{ocupado}(2,0)}A_2$,\ $A_1\xrightarrow{\mathit{livre}(1,1)}A_3$,\ $A_2\xrightarrow{\mathit{livre}(0,1)}A_3$,\ $A_3\xrightarrow{\mathit{ocupado}(4,1)}A_4$,\ $A_3\xrightarrow{\mathit{ocupado}(5,1)}A_5$,\ $A_4\xrightarrow{\mathit{livre}(2,1)}A_5$ \\
\midrule
Sit.\ 3 ($S_7$) & $A_1=\Move(d,c,0)$, $A_2=\Move(a,b,5)$, $A_3=\Move(d,T,2)$, $A_4=\Move(a,c,0)$, $A_5=\Move(b,c,1)$, $A_6=\Move(d,T,3)$ &
$A_1\xrightarrow{\mathit{livre}(3,1),\,\mathit{livre}(5,1)}A_2$,\ $A_1\xrightarrow{\mathit{livre}(3,1),\,\mathit{livre}(4,1)}A_6$,\ $A_2\xrightarrow{\mathit{livre}(3,0)}A_3$,\ $A_3\xrightarrow{\mathit{livre}(0,1)}A_4$,\ $A_3\xrightarrow{\mathit{livre}(1,1)}A_5$,\ $A_3\xrightarrow{\mathit{livre}(2,1)}A_6$,\ $A_4\xrightarrow{\mathit{livre}(5,1)}A_5$,\ $A_5\xrightarrow{\mathit{livre}(5,0)}A_6$ \\
\bottomrule
\end{tabular}
\end{center}

\paragraph{Resultado.} Nos quatro planos, as restrições de ordem necessárias formam uma cadeia $A_1\prec A_2\prec\dots\prec A_k$, e só existe uma ordem de execução válida das mesmas ações. Cada ação libera ou ocupa exatamente o espaço que a seguinte precisa, então não há ações independentes. Com os blocos e as metas dessas situações, o plano mínimo é uma cadeia. Os $A_i\prec A_j$ que não são consecutivos, como $A_1\prec A_6$ na Situação 3, já decorrem da transitividade da cadeia.

\subsubsection{Ordem parcial entre metas}

O enunciado pede que a solução considere ordem parcial, $A\xrightarrow{P}B$, e o manual do professor indica como codificá-la: ligar metas por $\varphi_1\prec_P\varphi_2$. Uma \emph{meta atômica} é $b(p,l)$, isto é, $\At(b,p)\land\Lev(b,l)$. Definimos
\[
\mathit{Ach}(\varphi,t)\leftrightarrow\exists t'\le t.\ \varphi(t')
\qquad
\varphi_1\prec_P\varphi_2\ \Longleftrightarrow\ \forall t.\ \bigl(\varphi_2(t)\rightarrow\mathit{Ach}(\varphi_1,t)\bigr)
\]
Em palavras, a meta $\varphi_2$ só pode valer num instante em que $\varphi_1$ já valeu antes ou naquele instante. Metas sem relação entre si ficam livres. A relação $\prec_P$ entre metas é irreflexiva e transitiva, como $\prec$ entre passos. Na prática, o plano deve alcançar $\varphi_1$ antes de alcançar $\varphi_2$.

\paragraph{Exemplo (Situação 3).} As metas são $c(0,0)$, $d(3,0)$, $a(0,1)$ e $b(1,1)$.
\begin{itemize}
\item Com $a(0,1)\prec_P b(1,1)$ (colocar $a$ antes de $b$), o plano mínimo continua com 6 ações, o mesmo do item 3, que já respeita essa ordem: $a(0,1)$ passa a valer em $t=4$ e $b(1,1)$ em $t=5$.
\item Com $b(1,1)\prec_P a(0,1)$ (colocar $b$ antes de $a$), não existe plano com 5 nem com 6 ações. O plano mínimo passa a ter 7: $b(1,1)$ passa a valer em $t=5$ e $a(0,1)$ em $t=6$. No plano encontrado, o bloco $a$ está sobre $b$ e precisa sair de cima dele antes de $b$ subir. Ele fica temporariamente sobre $d$ e só depois vai para $c$, o que custa uma ação a mais.
\end{itemize}
Os tamanhos 5, 6 e 7 vêm da execução do SAT solver para $T=5,6,7$ (seção~\ref{sec:resultados}); o plano de 7 ações foi conferido contra as pré-condições P1 a P6.

\paragraph{Plano de 7 ações como ordem parcial.} O plano com 7 ações é um caso em que os elos causais deixam ações livres. Para o plano devolvido pelo SAT, $A_1=\Move(d,c,0)$, $A_2=\Move(a,b,5)$, $A_3=\Move(d,T,2)$, $A_4=\Move(a,d,2)$, $A_5=\Move(b,c,1)$, $A_6=\Move(a,c,0)$ e $A_7=\Move(d,T,3)$, as restrições necessárias são
\[
A_1\prec A_2,\ A_1\prec A_4,\ A_2\prec A_3,\ A_3\prec A_5,\ A_3\prec A_6,\ A_4\prec A_5,\ A_4\prec A_6,\ A_5\prec A_7,\ A_6\prec A_7
\]
Existem 3 ordens de execução fisicamente válidas. Só uma delas, $A_1A_2A_3A_4A_5A_6A_7$, executa $A_5$ antes de $A_6$, como a meta $b(1,1)\prec_P a(0,1)$ exige. Assim, a ordem parcial entre metas funciona como uma restrição extra, $A_5\prec A_6$, que escolhe uma entre as ordens permitidas pelos elos causais.


% RESPONSAVEL: Guilherme (item 3)
\subsection{Execução Manual dos Cenários}

Usamos a representação do Item 1 e os efeitos do Item 2. Nível máximo 2 e mesa de 0 a 6; cada bloco é escrito como $b(p,l)$.
Em $\Poss(\Move(b,y,p),t)$, o nível-alvo é $L=0$ se $y=T$ e $L=l_y+1$ se $y$ é bloco. As pré-condições, numeradas, são:
\begin{enumerate}
\item[(P1)] $\Clr(b,t)$
\item[(P2)] $y\neq b\ \land\ 0\le p\le 6-\ell(b)\ \land\ L\le 2$
\item[(P3)] $y=T\ \lor\ \exists q.\ \At(y,q,t)\land\Overlap(b,p,y,q)$
\item[(P4)] $\forall s\in[p,p+\ell(b)).\ \Free(s,L,b,t)$
\item[(P5)] $\Stable(b,p,L,t)$
\item[(P6)] $\neg(\At(b,p,t)\land\Lev(b,L,t))$
\end{enumerate}
Efeitos de $t$ para $t+1$: ADD $\At(b,p,t+1)$ e $\Lev(b,L,t+1)$. DELETE da posição anterior $p_0$ e do nível anterior $l_0$, \emph{somente se forem diferentes} de $p$ e $L$ (se $p_0=p$, o $\At$ não é apagado; se $l_0=L$, o $\Lev$ não é apagado). Persistência: todo $b'\neq b$ mantém $\At$ e $\Lev$. Os predicados derivados são recalculados.

\subsubsection{Situação 1: $S_0\to S_{f3}$}
Estado inicial e final conforme o item 1 ($S_{f3}$ é o mais curto dos quatro; os outros levam 4, 8 e 9 movimentos).

Estado inicial: $a(3,0)\ b(5,0)\ c(0,0)\ d(3,1)$. Estado final: $a(2,0)\ b(5,0)\ c(0,0)\ d(0,1)$.

\begin{longtable}{c p{2.6cm} c p{3.4cm} p{3.4cm} p{3.6cm}}
\toprule $t$ & Ação & P1--P6 & Add & Delete & Estado em $t+1$ \\ \midrule \endhead
0 & $\Move(d,c,0)$ & $\checkmark$ & $\At(d,0,1),\ \Lev(d,1,1)$ & $\At(d,3,1)$ & $a(3,0)\ b(5,0)\ c(0,0)\ d(0,1)$ \\
1 & $\Move(a,T,2)$ & $\checkmark$ & $\At(a,2,2),\ \Lev(a,0,2)$ & $\At(a,3,2)$ & $a(2,0)\ b(5,0)\ c(0,0)\ d(0,1)$ \\
\bottomrule \end{longtable}
\paragraph{Verificação da $\Poss$.}\begin{itemize}
\item $t=0$: $\Move(d,c,0)$ ($L=1$; antes: $d(3,1)$).
\begin{itemize}
\item (P1) nada cobre os slots \{3,4,5\} no nível 2\ $\checkmark$
\item (P2) y$\neq$d, 0 $\le$ 0 $\le$ 3, L=1 $\le$ 2\ $\checkmark$
\item (P3) span(d,0) = [0,3) e span(c,0) = [0,2) se sobrepõem\ $\checkmark$
\item (P4) slots \{0,1,2\} livres no nível 1\ $\checkmark$
\item (P5) slots \{0,1\} apoiados: 2 $\ge$ $\lceil$3/2$\rceil$ = 2\ $\checkmark$
\item (P6) (3,1) $\neq$ (0,1)\ $\checkmark$
\end{itemize}
\item $t=1$: $\Move(a,T,2)$ ($L=0$; antes: $a(3,0)$).
\begin{itemize}
\item (P1) nada cobre os slots \{3\} no nível 1\ $\checkmark$
\item (P2) y$\neq$a, 0 $\le$ 2 $\le$ 5, L=0 $\le$ 2\ $\checkmark$
\item (P3) y = T\ $\checkmark$
\item (P4) slots \{2\} livres no nível 0\ $\checkmark$
\item (P5) L = 0 (apoio na mesa)\ $\checkmark$
\item (P6) (3,0) $\neq$ (2,0)\ $\checkmark$
\end{itemize}
\end{itemize}
Relações $\On$ no estado final: $\On(a,T), \On(b,T), \On(c,T), \On(d,a), \On(d,c)$.

\subsubsection{Situação 1 (extra): $S_0\to S_{f4}$}
Plano de 4 movimentos. Coincide com os estados $S_0$ a $S_4$ da Situação 3.

Estado inicial: $a(3,0)\ b(5,0)\ c(0,0)\ d(3,1)$. Estado final: $a(0,1)\ b(5,0)\ c(0,0)\ d(2,0)$.

\begin{longtable}{c p{2.6cm} c p{3.4cm} p{3.4cm} p{3.6cm}}
\toprule $t$ & Ação & P1--P6 & Add & Delete & Estado em $t+1$ \\ \midrule \endhead
0 & $\Move(d,c,0)$ & $\checkmark$ & $\At(d,0,1),\ \Lev(d,1,1)$ & $\At(d,3,1)$ & $a(3,0)\ b(5,0)\ c(0,0)\ d(0,1)$ \\
1 & $\Move(a,b,5)$ & $\checkmark$ & $\At(a,5,2),\ \Lev(a,1,2)$ & $\At(a,3,2),\ \Lev(a,0,2)$ & $a(5,1)\ b(5,0)\ c(0,0)\ d(0,1)$ \\
2 & $\Move(d,T,2)$ & $\checkmark$ & $\At(d,2,3),\ \Lev(d,0,3)$ & $\At(d,0,3),\ \Lev(d,1,3)$ & $a(5,1)\ b(5,0)\ c(0,0)\ d(2,0)$ \\
3 & $\Move(a,c,0)$ & $\checkmark$ & $\At(a,0,4),\ \Lev(a,1,4)$ & $\At(a,5,4)$ & $a(0,1)\ b(5,0)\ c(0,0)\ d(2,0)$ \\
\bottomrule \end{longtable}
\paragraph{Verificação da $\Poss$.}\begin{itemize}
\item $t=0$: $\Move(d,c,0)$ ($L=1$; antes: $d(3,1)$).
\begin{itemize}
\item (P1) nada cobre os slots \{3,4,5\} no nível 2\ $\checkmark$
\item (P2) y$\neq$d, 0 $\le$ 0 $\le$ 3, L=1 $\le$ 2\ $\checkmark$
\item (P3) span(d,0) = [0,3) e span(c,0) = [0,2) se sobrepõem\ $\checkmark$
\item (P4) slots \{0,1,2\} livres no nível 1\ $\checkmark$
\item (P5) slots \{0,1\} apoiados: 2 $\ge$ $\lceil$3/2$\rceil$ = 2\ $\checkmark$
\item (P6) (3,1) $\neq$ (0,1)\ $\checkmark$
\end{itemize}
\item $t=1$: $\Move(a,b,5)$ ($L=1$; antes: $a(3,0)$).
\begin{itemize}
\item (P1) nada cobre os slots \{3\} no nível 1\ $\checkmark$
\item (P2) y$\neq$a, 0 $\le$ 5 $\le$ 5, L=1 $\le$ 2\ $\checkmark$
\item (P3) span(a,5) = [5,6) e span(b,5) = [5,6) se sobrepõem\ $\checkmark$
\item (P4) slots \{5\} livres no nível 1\ $\checkmark$
\item (P5) slots \{5\} apoiados: 1 $\ge$ $\lceil$1/2$\rceil$ = 1\ $\checkmark$
\item (P6) (3,0) $\neq$ (5,1)\ $\checkmark$
\end{itemize}
\item $t=2$: $\Move(d,T,2)$ ($L=0$; antes: $d(0,1)$).
\begin{itemize}
\item (P1) nada cobre os slots \{0,1,2\} no nível 2\ $\checkmark$
\item (P2) y$\neq$d, 0 $\le$ 2 $\le$ 3, L=0 $\le$ 2\ $\checkmark$
\item (P3) y = T\ $\checkmark$
\item (P4) slots \{2,3,4\} livres no nível 0\ $\checkmark$
\item (P5) L = 0 (apoio na mesa)\ $\checkmark$
\item (P6) (0,1) $\neq$ (2,0)\ $\checkmark$
\end{itemize}
\item $t=3$: $\Move(a,c,0)$ ($L=1$; antes: $a(5,1)$).
\begin{itemize}
\item (P1) nada cobre os slots \{5\} no nível 2\ $\checkmark$
\item (P2) y$\neq$a, 0 $\le$ 0 $\le$ 5, L=1 $\le$ 2\ $\checkmark$
\item (P3) span(a,0) = [0,1) e span(c,0) = [0,2) se sobrepõem\ $\checkmark$
\item (P4) slots \{0\} livres no nível 1\ $\checkmark$
\item (P5) slots \{0\} apoiados: 1 $\ge$ $\lceil$1/2$\rceil$ = 1\ $\checkmark$
\item (P6) (5,1) $\neq$ (0,1)\ $\checkmark$
\end{itemize}
\end{itemize}
Relações $\On$ no estado final: $\On(a,c), \On(b,T), \On(c,T), \On(d,T)$.

\subsubsection{Situação 2: $S_0\to S_5$}
Plano de 5 movimentos, que passa pelos estados $S_1$ a $S_4$ do slide.

Estado inicial: $a(0,1)\ b(1,1)\ c(0,0)\ d(3,0)$. Estado final: $a(4,2)\ b(5,2)\ c(4,1)\ d(3,0)$.

\begin{longtable}{c p{2.6cm} c p{3.4cm} p{3.4cm} p{3.6cm}}
\toprule $t$ & Ação & P1--P6 & Add & Delete & Estado em $t+1$ \\ \midrule \endhead
0 & $\Move(b,T,2)$ & $\checkmark$ & $\At(b,2,1),\ \Lev(b,0,1)$ & $\At(b,1,1),\ \Lev(b,1,1)$ & $a(0,1)\ b(2,0)\ c(0,0)\ d(3,0)$ \\
1 & $\Move(a,b,2)$ & $\checkmark$ & $\At(a,2,2),\ \Lev(a,1,2)$ & $\At(a,0,2)$ & $a(2,1)\ b(2,0)\ c(0,0)\ d(3,0)$ \\
2 & $\Move(c,d,4)$ & $\checkmark$ & $\At(c,4,3),\ \Lev(c,1,3)$ & $\At(c,0,3),\ \Lev(c,0,3)$ & $a(2,1)\ b(2,0)\ c(4,1)\ d(3,0)$ \\
3 & $\Move(a,c,4)$ & $\checkmark$ & $\At(a,4,4),\ \Lev(a,2,4)$ & $\At(a,2,4),\ \Lev(a,1,4)$ & $a(4,2)\ b(2,0)\ c(4,1)\ d(3,0)$ \\
4 & $\Move(b,c,5)$ & $\checkmark$ & $\At(b,5,5),\ \Lev(b,2,5)$ & $\At(b,2,5),\ \Lev(b,0,5)$ & $a(4,2)\ b(5,2)\ c(4,1)\ d(3,0)$ \\
\bottomrule \end{longtable}
\paragraph{Verificação da $\Poss$.}\begin{itemize}
\item $t=0$: $\Move(b,T,2)$ ($L=0$; antes: $b(1,1)$).
\begin{itemize}
\item (P1) nada cobre os slots \{1\} no nível 2\ $\checkmark$
\item (P2) y$\neq$b, 0 $\le$ 2 $\le$ 5, L=0 $\le$ 2\ $\checkmark$
\item (P3) y = T\ $\checkmark$
\item (P4) slots \{2\} livres no nível 0\ $\checkmark$
\item (P5) L = 0 (apoio na mesa)\ $\checkmark$
\item (P6) (1,1) $\neq$ (2,0)\ $\checkmark$
\end{itemize}
\item $t=1$: $\Move(a,b,2)$ ($L=1$; antes: $a(0,1)$).
\begin{itemize}
\item (P1) nada cobre os slots \{0\} no nível 2\ $\checkmark$
\item (P2) y$\neq$a, 0 $\le$ 2 $\le$ 5, L=1 $\le$ 2\ $\checkmark$
\item (P3) span(a,2) = [2,3) e span(b,2) = [2,3) se sobrepõem\ $\checkmark$
\item (P4) slots \{2\} livres no nível 1\ $\checkmark$
\item (P5) slots \{2\} apoiados: 1 $\ge$ $\lceil$1/2$\rceil$ = 1\ $\checkmark$
\item (P6) (0,1) $\neq$ (2,1)\ $\checkmark$
\end{itemize}
\item $t=2$: $\Move(c,d,4)$ ($L=1$; antes: $c(0,0)$).
\begin{itemize}
\item (P1) nada cobre os slots \{0,1\} no nível 1\ $\checkmark$
\item (P2) y$\neq$c, 0 $\le$ 4 $\le$ 4, L=1 $\le$ 2\ $\checkmark$
\item (P3) span(c,4) = [4,6) e span(d,3) = [3,6) se sobrepõem\ $\checkmark$
\item (P4) slots \{4,5\} livres no nível 1\ $\checkmark$
\item (P5) slots \{4,5\} apoiados: 2 $\ge$ $\lceil$2/2$\rceil$ = 1\ $\checkmark$
\item (P6) (0,0) $\neq$ (4,1)\ $\checkmark$
\end{itemize}
\item $t=3$: $\Move(a,c,4)$ ($L=2$; antes: $a(2,1)$).
\begin{itemize}
\item (P1) nada cobre os slots \{2\} no nível 2\ $\checkmark$
\item (P2) y$\neq$a, 0 $\le$ 4 $\le$ 5, L=2 $\le$ 2\ $\checkmark$
\item (P3) span(a,4) = [4,5) e span(c,4) = [4,6) se sobrepõem\ $\checkmark$
\item (P4) slots \{4\} livres no nível 2\ $\checkmark$
\item (P5) slots \{4\} apoiados: 1 $\ge$ $\lceil$1/2$\rceil$ = 1\ $\checkmark$
\item (P6) (2,1) $\neq$ (4,2)\ $\checkmark$
\end{itemize}
\item $t=4$: $\Move(b,c,5)$ ($L=2$; antes: $b(2,0)$).
\begin{itemize}
\item (P1) nada cobre os slots \{2\} no nível 1\ $\checkmark$
\item (P2) y$\neq$b, 0 $\le$ 5 $\le$ 5, L=2 $\le$ 2\ $\checkmark$
\item (P3) span(b,5) = [5,6) e span(c,4) = [4,6) se sobrepõem\ $\checkmark$
\item (P4) slots \{5\} livres no nível 2\ $\checkmark$
\item (P5) slots \{5\} apoiados: 1 $\ge$ $\lceil$1/2$\rceil$ = 1\ $\checkmark$
\item (P6) (2,0) $\neq$ (5,2)\ $\checkmark$
\end{itemize}
\end{itemize}
Relações $\On$ no estado final: $\On(a,c), \On(b,c), \On(c,d), \On(d,T)$.

\subsubsection{Situação 3: $S_0\to S_7$}
Plano de 6 movimentos: $S_0$$\to$$S_1$$\to$$S_2$$\to$$S_3$$\to$$S_4$$\to$$S_5$$\to$$S_7$ ($S_6$ é igual a $S_5$, sem ação entre eles).

Estado inicial: $a(3,0)\ b(5,0)\ c(0,0)\ d(3,1)$. Estado final: $a(0,1)\ b(1,1)\ c(0,0)\ d(3,0)$.

\begin{longtable}{c p{2.6cm} c p{3.4cm} p{3.4cm} p{3.6cm}}
\toprule $t$ & Ação & P1--P6 & Add & Delete & Estado em $t+1$ \\ \midrule \endhead
0 & $\Move(d,c,0)$ & $\checkmark$ & $\At(d,0,1),\ \Lev(d,1,1)$ & $\At(d,3,1)$ & $a(3,0)\ b(5,0)\ c(0,0)\ d(0,1)$ \\
1 & $\Move(a,b,5)$ & $\checkmark$ & $\At(a,5,2),\ \Lev(a,1,2)$ & $\At(a,3,2),\ \Lev(a,0,2)$ & $a(5,1)\ b(5,0)\ c(0,0)\ d(0,1)$ \\
2 & $\Move(d,T,2)$ & $\checkmark$ & $\At(d,2,3),\ \Lev(d,0,3)$ & $\At(d,0,3),\ \Lev(d,1,3)$ & $a(5,1)\ b(5,0)\ c(0,0)\ d(2,0)$ \\
3 & $\Move(a,c,0)$ & $\checkmark$ & $\At(a,0,4),\ \Lev(a,1,4)$ & $\At(a,5,4)$ & $a(0,1)\ b(5,0)\ c(0,0)\ d(2,0)$ \\
4 & $\Move(b,c,1)$ & $\checkmark$ & $\At(b,1,5),\ \Lev(b,1,5)$ & $\At(b,5,5),\ \Lev(b,0,5)$ & $a(0,1)\ b(1,1)\ c(0,0)\ d(2,0)$ \\
5 & $\Move(d,T,3)$ & $\checkmark$ & $\At(d,3,6),\ \Lev(d,0,6)$ & $\At(d,2,6)$ & $a(0,1)\ b(1,1)\ c(0,0)\ d(3,0)$ \\
\bottomrule \end{longtable}
\paragraph{Verificação da $\Poss$.}\begin{itemize}
\item $t=0$: $\Move(d,c,0)$ ($L=1$; antes: $d(3,1)$).
\begin{itemize}
\item (P1) nada cobre os slots \{3,4,5\} no nível 2\ $\checkmark$
\item (P2) y$\neq$d, 0 $\le$ 0 $\le$ 3, L=1 $\le$ 2\ $\checkmark$
\item (P3) span(d,0) = [0,3) e span(c,0) = [0,2) se sobrepõem\ $\checkmark$
\item (P4) slots \{0,1,2\} livres no nível 1\ $\checkmark$
\item (P5) slots \{0,1\} apoiados: 2 $\ge$ $\lceil$3/2$\rceil$ = 2\ $\checkmark$
\item (P6) (3,1) $\neq$ (0,1)\ $\checkmark$
\end{itemize}
\item $t=1$: $\Move(a,b,5)$ ($L=1$; antes: $a(3,0)$).
\begin{itemize}
\item (P1) nada cobre os slots \{3\} no nível 1\ $\checkmark$
\item (P2) y$\neq$a, 0 $\le$ 5 $\le$ 5, L=1 $\le$ 2\ $\checkmark$
\item (P3) span(a,5) = [5,6) e span(b,5) = [5,6) se sobrepõem\ $\checkmark$
\item (P4) slots \{5\} livres no nível 1\ $\checkmark$
\item (P5) slots \{5\} apoiados: 1 $\ge$ $\lceil$1/2$\rceil$ = 1\ $\checkmark$
\item (P6) (3,0) $\neq$ (5,1)\ $\checkmark$
\end{itemize}
\item $t=2$: $\Move(d,T,2)$ ($L=0$; antes: $d(0,1)$).
\begin{itemize}
\item (P1) nada cobre os slots \{0,1,2\} no nível 2\ $\checkmark$
\item (P2) y$\neq$d, 0 $\le$ 2 $\le$ 3, L=0 $\le$ 2\ $\checkmark$
\item (P3) y = T\ $\checkmark$
\item (P4) slots \{2,3,4\} livres no nível 0\ $\checkmark$
\item (P5) L = 0 (apoio na mesa)\ $\checkmark$
\item (P6) (0,1) $\neq$ (2,0)\ $\checkmark$
\end{itemize}
\item $t=3$: $\Move(a,c,0)$ ($L=1$; antes: $a(5,1)$).
\begin{itemize}
\item (P1) nada cobre os slots \{5\} no nível 2\ $\checkmark$
\item (P2) y$\neq$a, 0 $\le$ 0 $\le$ 5, L=1 $\le$ 2\ $\checkmark$
\item (P3) span(a,0) = [0,1) e span(c,0) = [0,2) se sobrepõem\ $\checkmark$
\item (P4) slots \{0\} livres no nível 1\ $\checkmark$
\item (P5) slots \{0\} apoiados: 1 $\ge$ $\lceil$1/2$\rceil$ = 1\ $\checkmark$
\item (P6) (5,1) $\neq$ (0,1)\ $\checkmark$
\end{itemize}
\item $t=4$: $\Move(b,c,1)$ ($L=1$; antes: $b(5,0)$).
\begin{itemize}
\item (P1) nada cobre os slots \{5\} no nível 1\ $\checkmark$
\item (P2) y$\neq$b, 0 $\le$ 1 $\le$ 5, L=1 $\le$ 2\ $\checkmark$
\item (P3) span(b,1) = [1,2) e span(c,0) = [0,2) se sobrepõem\ $\checkmark$
\item (P4) slots \{1\} livres no nível 1\ $\checkmark$
\item (P5) slots \{1\} apoiados: 1 $\ge$ $\lceil$1/2$\rceil$ = 1\ $\checkmark$
\item (P6) (5,0) $\neq$ (1,1)\ $\checkmark$
\end{itemize}
\item $t=5$: $\Move(d,T,3)$ ($L=0$; antes: $d(2,0)$).
\begin{itemize}
\item (P1) nada cobre os slots \{2,3,4\} no nível 1\ $\checkmark$
\item (P2) y$\neq$d, 0 $\le$ 3 $\le$ 3, L=0 $\le$ 2\ $\checkmark$
\item (P3) y = T\ $\checkmark$
\item (P4) slots \{3,4,5\} livres no nível 0\ $\checkmark$
\item (P5) L = 0 (apoio na mesa)\ $\checkmark$
\item (P6) (2,0) $\neq$ (3,0)\ $\checkmark$
\end{itemize}
\end{itemize}
Relações $\On$ no estado final: $\On(a,c), \On(b,c), \On(c,T), \On(d,T)$.

% RASCUNHO da gerencia (Guilherme) para o item 5 -- revisar com Yuri e Gabriel.
% Descreve a codificacao feita por gerar_cnf.py (versao com as correcoes de estabilidade e de nivel maximo).
\section{Codificação CNF}

Para um horizonte de $T$ passos, o problema de planejamento vira um problema de satisfatibilidade:
existem os instantes $0,1,\dots,T$ e as ações ocorrem nos instantes $0,\dots,T-1$.
A fórmula é satisfatível se e somente se existe um plano com exatamente $T$ ações que leva do estado inicial ao estado meta.
Cada ação $\Move(b,y,p)$ do item 2 vira uma variável $\mathit{mv}(b,y,p,t)$, e cada predicado da representação em LPO vira uma família de variáveis proposicionais indexadas pelo instante.

\subsection{Variáveis proposicionais}

Sejam $B=\{a,b,c,d\}$, os pontos $0,\dots,6$, os slots $0,\dots,5$ e os níveis $0,1,2$.
\begin{center}
\begin{tabular}{p{3.6cm} p{6.6cm} p{3.4cm}}
\toprule
Variável & Índices & Quantidade \\
\midrule
$\At(b,p,t)$ & $b\in B$, $p\in[0,6-\ell(b)]$, $t\in[0,T]$ & $21\,(T+1)$ \\
$\Lev(b,l,t)$ & $b\in B$, $l\in[0,2]$, $t\in[0,T]$ & $12\,(T+1)$ \\
$\Clr(b,t)$ & $b\in B$, $t\in[0,T]$ & $4\,(T+1)$ \\
$\Cov(b,s,l,t)$ & $b\in B$, $s\in[0,5]$, $l\in[0,2]$, $t\in[0,T]$ & $72\,(T+1)$ \\
$\mathit{mv}(b,y,p,t)$ & $b\in B$, $y\in (B\setminus\{b\})\cup\{T\}$, $p\in[0,6-\ell(b)]$, $t\in[0,T-1]$ & $84\,T$ \\
\midrule
\multicolumn{2}{l}{Total} & $193\,T+109$ \\
\bottomrule
\end{tabular}
\end{center}
Para $T=2$, $T=5$ e $T=6$ isso dá 495, 1074 e 1267 variáveis.

Os predicados $\Free$, $\Stable$, $\Span$ e $\Overlap$ não viram variáveis: o gerador os expande diretamente em cláusulas. A relação $\On$ também não é codificada (seção~\ref{sec:on}).
A variável $\Cov$ foi acrescentada em relação à tabela do manual do professor: ela permite exigir que o apoio de um bloco esteja no nível imediatamente abaixo dele, e não em qualquer nível.
Quando há ordem parcial entre metas (grupo 11), entram duas variáveis por meta atômica e por instante, $\mathit{hold}(b,p,l,t)$ e $\mathit{ach}(b,p,l,t)$, isto é, $2(T+1)$ por meta.

\subsection{Grupos de cláusulas}

Em todos os grupos, $t$ percorre os instantes em que a regra vale. Os grupos 1 a 6 descrevem os estados, os grupos 7 a 10 descrevem as ações e o grupo 11, opcional, trata a ordem parcial entre metas.

\begin{enumerate}
\item \textbf{Estado inicial.} Para cada bloco $b$ com $b(p_0,l_0)$ no estado inicial, as cláusulas unitárias $\At(b,p_0,0)$ e $\Lev(b,l_0,0)$.
\item \textbf{Meta.} Para cada bloco $b$ com $b(p_f,l_f)$ no estado meta, as cláusulas unitárias $\At(b,p_f,T)$ e $\Lev(b,l_f,T)$.
\item \textbf{Unicidade de posição e de nível} (leis U1 e U2). Para cada $b$ e $t$:
\[
\bigvee_{p}\At(b,p,t)\qquad \neg\At(b,p_1,t)\lor\neg\At(b,p_2,t)\ (p_1<p_2)
\]
e o análogo para $\Lev$ com $l\in\{0,1,2\}$.
\item \textbf{Definição de $\Cov$.} Para cada $b$, $s$, $l$ e $t$, com $P_s=\{p : p\le s<p+\ell(b)\}$:
\begin{align*}
&\Cov(b,s,l,t)\rightarrow\Lev(b,l,t)\\
&\Cov(b,s,l,t)\rightarrow\bigvee_{p\in P_s}\At(b,p,t)\\
&\At(b,p,t)\land\Lev(b,l,t)\rightarrow\Cov(b,s,l,t)\qquad(p\in P_s)
\end{align*}
\item \textbf{Exclusão horizontal} (lei E). Para blocos distintos $b,b'$, posições $p,q$ com $\Overlap(b,p,b',q)$ e todo nível $l$:
\[
\neg\At(b,p,t)\lor\neg\Lev(b,l,t)\lor\neg\At(b',q,t)\lor\neg\Lev(b',l,t)
\]
\item \textbf{Clear.} Se $\Clr(b,t)$ vale, nenhum outro bloco $z$ está no nível seguinte sobreposto a $b$. Para $\Overlap(b,p,z,q)$ e $l<2$:
\begin{align*}
\neg\Clr(b,t)\lor\neg\At(b,p,t)\lor\neg\Lev(b,l,t)\lor\neg\At(z,q,t)\lor\neg\Lev(z,l+1,t)
\end{align*}
O gerador escreve também cláusulas na direção contrária, mas elas não são necessárias para a correção: $\Clr$ só aparece como pré-condição de $\Move$, e o solver pode escolher $\Clr(b,t)$ verdadeiro sempre que nada está sobre $b$.
\item \textbf{Ação única por instante.} Para cada $t\in[0,T-1]$, pelo menos uma variável $\mathit{mv}(\cdot,t)$ é verdadeira, e para cada par de ações distintas $\neg\mathit{mv}_1\lor\neg\mathit{mv}_2$.
\item \textbf{Pré-condições de $\Move$.} Para cada $\mathit{mv}(b,y,p,t)$, com $L=0$ se $y=T$ e $L=l_y+1$ caso contrário:
\begin{itemize}
\item[(P1)] $\mathit{mv}(b,y,p,t)\rightarrow\Clr(b,t)$.
\item[(P2)] $y\neq b$ e $0\le p\le 6-\ell(b)$ já valem pelo domínio das variáveis. O limite $L\le 2$ vira $\mathit{mv}(b,y,p,t)\rightarrow\neg\Lev(y,2,t)$.
\item[(P3)] Se $y\neq T$, para cada posição $q$ de $y$ sem sobreposição: $\neg\mathit{mv}(b,y,p,t)\lor\neg\Lev(y,l_y,t)\lor\neg\At(y,q,t)$.
\item[(P4)] Para cada outro bloco $z$ e posição $q$ com $\Overlap(b,p,z,q)$: $\neg\mathit{mv}(b,y,p,t)\lor\neg\Lev(y,l_y,t)\lor\neg\At(z,q,t)\lor\neg\Lev(z,L,t)$. Para $y=T$, o mesmo com $L=0$ e sem a parcela de $y$.
\item[(P5)] Estabilidade, descrita a seguir.
\item[(P6)] $\neg\mathit{mv}(b,y,p,t)\lor\neg\At(b,p,t)\lor\neg\Lev(y,l_y,t)\lor\neg\Lev(b,L,t)$, o que evita ações que não mudam o estado.
\end{itemize}
\item \textbf{Efeitos de $\Move$.} $\mathit{mv}(b,y,p,t)\rightarrow\At(b,p,t+1)$, e $\mathit{mv}(b,y,p,t)\land\Lev(y,l_y,t)\rightarrow\Lev(b,l_y+1,t+1)$. Para $y=T$, $\mathit{mv}(b,T,p,t)\rightarrow\Lev(b,0,t+1)$.
\item \textbf{Frame axioms.} Para $b'\neq b$, $\mathit{mv}(b,y,p,t)\land\At(b',p',t)\rightarrow\At(b',p',t+1)$ e $\mathit{mv}(b,y,p,t)\land\Lev(b',l',t)\rightarrow\Lev(b',l',t+1)$.
\item \textbf{Ordem parcial entre metas (opcional).} Para cada relação $\varphi_1\prec_P\varphi_2$, com metas atômicas $\varphi=b(p,l)$, cláusulas que dizem que $\varphi_2$ só vale em $t$ se $\varphi_1$ já valeu em algum $t'\le t$:
\begin{align*}
&\mathit{hold}(\varphi,t)\rightarrow\At(b,p,t) \qquad \mathit{hold}(\varphi,t)\rightarrow\Lev(b,l,t) \qquad \At(b,p,t)\land\Lev(b,l,t)\rightarrow\mathit{hold}(\varphi,t)\\
&\mathit{ach}(\varphi,0)\rightarrow\mathit{hold}(\varphi,0) \qquad \mathit{ach}(\varphi,t)\rightarrow\mathit{ach}(\varphi,t-1)\lor\mathit{hold}(\varphi,t)\ (t\ge 1)\\
&\mathit{hold}(\varphi_2,t)\rightarrow\mathit{ach}(\varphi_1,t)\qquad(t\in[0,T])
\end{align*}
Basta a direção ``$\mathit{ach}$ implica que já valeu'': como $\mathit{ach}$ só aparece positivo na última cláusula, o solver só consegue torná-lo verdadeiro se $\varphi_1$ de fato valeu antes. É o que o manual do professor descreve como $\neg\varphi_2(t)\lor\varphi_1(t')$ para algum $t'\le t$.
\end{enumerate}

\subsection{Estabilidade em cláusulas}

A condição (P5) exige que pelo menos $k=\lceil\ell(b)/2\rceil$ dos $n=\ell(b)$ slots sob $b$ estejam apoiados no nível $l_y$.
``Pelo menos $k$ de $n$ slots apoiados'' equivale a ``todo subconjunto de $n-k+1$ slots tem algum slot apoiado''. Então, para cada subconjunto $S$ de $n-k+1$ slots de $\Span(b,p)$:
\[
\neg\mathit{mv}(b,y,p,t)\lor\neg\Lev(y,l_y,t)\lor\bigvee_{s\in S}\ \bigvee_{z\neq b}\Cov(z,s,l_y,t)
\]
Para $b=d$ ($n=3$, $k=2$) são três cláusulas, uma para cada par de slots. Para $a$, $b$ e $c$ ($k=1$ e $n-k+1=n$) é uma cláusula só, e ela exige apoio em algum dos slots sob o bloco.
Quando o destino é a mesa ($y=T$), não há cláusula de estabilidade.

\subsection{Como ADD e DELETE aparecem na codificação}

O ADD do item 2 são as cláusulas de efeito do grupo 9. Não existe cláusula de DELETE explícita: a posição e o nível anteriores deixam de valer porque as cláusulas de unicidade (grupo 3) proíbem duas posições ou dois níveis para o mesmo bloco no mesmo instante.
Isso é coerente com a regra do item 2 de só apagar o que muda: se $p_0=p$, o $\At$ novo e o antigo são o mesmo átomo e continuam verdadeiros.

\subsection{A relação $\On$ é derivada}\label{sec:on}

Não criamos variáveis para $\On(b,y,t)$. A relação é reconstruída depois, a partir do modelo do SAT solver:
\[
\On(b,y,t)\leftrightarrow\exists s,l.\ \Cov(b,s,l,t)\land\Cov(y,s,l-1,t)\qquad
\On(b,T,t)\leftrightarrow\Lev(b,0,t)
\]
Isso mantém o número de variáveis menor, e o script \texttt{interpretar.py} faz a reconstrução ao apresentar o plano.

\subsection{Horizonte e plano mínimo}

Como cada instante tem exatamente uma ação, a fórmula para $T$ é satisfatível quando existe plano de exatamente $T$ ações.
Para achar o plano mais curto, executa-se o SAT solver com $T=1,2,3,\dots$ e para-se no primeiro $T$ satisfatível. Esse $T$ é o comprimento mínimo do plano, porque os valores menores deram insatisfatível.

% RASCUNHO da gerencia (Guilherme) para o item 5 -- revisar com Yuri e Gabriel.
% Contagens obtidas com gerar_cnf.py (versao com cov e ordem parcial).
\section{Exemplos dos 3 Cenários Codificados para CNF}

Cada cenário vira um arquivo \texttt{.cnf} (as cláusulas) e um \texttt{.map} (o nome de cada variável). O estado inicial e o estado meta entram como cláusulas unitárias nos instantes $0$ e $T$, e as demais cláusulas são as do grupo~3 em diante da seção anterior. O horizonte $T$ usado em cada caso é o comprimento do plano mínimo.

\begin{center}
\begin{tabular}{l l l c r r}
\toprule
Cenário & Estado inicial & Estado meta & $T$ & Variáveis & Cláusulas \\
\midrule
Situação 1 ($S_0\to S_{f3}$) & $c(0,0)\ a(3,0)\ b(5,0)\ d(3,1)$ & $c(0,0)\ a(2,0)\ b(5,0)\ d(0,1)$ & 2 & 495 & 17\,648 \\
Situação 2 ($S_0\to S_5$) & $c(0,0)\ a(0,1)\ b(1,1)\ d(3,0)$ & $c(4,1)\ a(4,2)\ b(5,2)\ d(3,0)$ & 5 & 1\,074 & 42\,893 \\
Situação 3 ($S_0\to S_7$) & $c(0,0)\ a(3,0)\ b(5,0)\ d(3,1)$ & $c(0,0)\ a(0,1)\ b(1,1)\ d(3,0)$ & 6 & 1\,267 & 51\,308 \\
\bottomrule
\end{tabular}
\end{center}

\subsection{Situação 1}
Estado inicial $S_0$: $c(0,0)\ a(3,0)\ b(5,0)\ d(3,1)$, com $d$ em ponte sobre $a$ e $b$. Meta $S_{f3}$: $c(0,0)\ a(2,0)\ b(5,0)\ d(0,1)$. Horizonte $T=2$.
O arquivo \texttt{trab01\_blocos2SAT\_situacao1.map} associa, por exemplo, \texttt{480} a \lstinline|move(d,c,0,0)| e \texttt{369} a \lstinline|move(a,T,2,1)|.
As 16 cláusulas unitárias são as do estado inicial (8, em $t=0$) e as da meta (8, em $t=2$):
\begin{lstlisting}
t=0:  at(a,3,0) lev(a,0,0)  at(b,5,0) lev(b,0,0)  at(c,0,0) lev(c,0,0)  at(d,3,0) lev(d,1,0)
t=2:  at(a,2,2) lev(a,0,2)  at(b,5,2) lev(b,0,2)  at(c,0,2) lev(c,0,2)  at(d,0,2) lev(d,1,2)
\end{lstlisting}
As outras três metas da Situação 1 ($S_{f1}$, $S_{f2}$ e $S_{f4}$) são geradas com o terceiro argumento, por exemplo \texttt{python3 gerar\_cnf.py 1 4 f4}.

\subsection{Situação 2}
Estado inicial $S_0$: $c(0,0)\ a(0,1)\ b(1,1)\ d(3,0)$, com $a$ e $b$ lado a lado sobre $c$. Meta $S_5$: $c(4,1)\ a(4,2)\ b(5,2)\ d(3,0)$, com $c$ sobre $d$ e $a$ e $b$ lado a lado sobre $c$. Horizonte $T=5$.

\subsection{Situação 3}
Estado inicial $S_0$: o mesmo da Situação 1. Meta $S_7$: $c(0,0)\ a(0,1)\ b(1,1)\ d(3,0)$, com $a$ e $b$ lado a lado sobre $c$ e $d$ na mesa. Horizonte $T=6$. O estado $S_6$ do slide é igual a $S_5$ e não gera ação.

\subsection{Cenário com ordem parcial entre metas}
O grupo (ORD) acrescenta variáveis e cláusulas só quando há ordem entre metas. Para a Situação 3 com $b(1,1)\prec_P a(0,1)$ e $T=7$, o CNF tem 1\,492 variáveis e 59\,795 cláusulas, contra 1\,267 e 51\,308 do caso sem ordem com $T=6$. O comando é \texttt{python3 gerar\_cnf.py 3 7 --ordem='b(1,1)<a(0,1)'}.

% RASCUNHO da gerencia (Guilherme) para o item 5 -- revisar com Yuri e Gabriel.
% Os trechos de codigo sao identificados pelos comentarios de gerar_cnf.py, e nao por numero de linha.
\section{Mapeamento: Descrição Formal $\to$ Código}

A tabela liga cada elemento da descrição formal (itens 1 e 2) ao trecho correspondente de \texttt{gerar\_cnf.py}. Os blocos do código são identificados pelos comentários do próprio arquivo.

\begin{longtable}{p{5.2cm} p{9.6cm}}
\toprule
Regra formal & Código Python (\texttt{gerar\_cnf.py}) \\
\midrule \endhead
Blocos e comprimentos $\ell(b)$ & \lstinline|BLOCKS = {"a":1, "b":1, "c":2, "d":3}| \\
Pontos do eixo e nível máximo & \lstinline|MAX_POINT = 6|, \lstinline|MAX_LEVEL = 2| \\
Posições válidas $p\in[0,6-\ell(b)]$ & \lstinline|positions(b)| \\
$\Overlap(b,p,y,q)$ & \lstinline|overlap(b, p, y, q)| \\
$\At$, $\Lev$, $\Clr$, $\Cov$, $\mathit{mv}$ & dicionários \lstinline|at|, \lstinline|lev|, \lstinline|clr|, \lstinline|cov|, \lstinline|mv|, criados com \lstinline|C.new_var(...)| \\
Estado inicial e meta (grupos 1 e 2) & bloco \lstinline|# Estado inicial e estado meta| \\
Unicidade de posição e nível, U1 e U2 (grupo 3) & bloco \lstinline|# (A,B)| \\
Definição de $\Cov$ (grupo 4) & bloco \lstinline|# (COV)| \\
Exclusão horizontal, lei E (grupo 5) & bloco \lstinline|# (C)| \\
$\Clr$ (grupo 6) & bloco \lstinline|# (E)| \\
Ação única por instante (grupo 7) & início do laço \texttt{\# Uma a\c{c}\~ao por instante...} \\
P1: $\Clr(b,t)$ & \lstinline|C.add(-v, clr[b, t])| \\
P2: $L\le 2$ & \lstinline|C.add(-v, -lev[y, MAX_LEVEL, t])| \\
P3: sobreposição com o apoio & laço \lstinline|for q in positions(y): if not overlap(...)| \\
P4: $\Free$ nos slots de destino & trechos \lstinline|# Destino na mesa...| e \texttt{\# N\~ao coincidir horizontalmente...} \\
P5: $\Stable$ & bloco \lstinline|# Estabilidade...|, com \lstinline|itertools.combinations(slots, r)| e \lstinline|cov[z, s, ly, t]| \\
P6: $\neg(\At\land\Lev)$ & bloco \lstinline|# Evita no-op...| \\
Efeitos ADD de $\Move$ (grupo 9) & \lstinline|C.add(-v, at[b, p, t+1])| e \lstinline|C.add(-v, lev[b, ..., t+1])| \\
DELETE & não há código: vem da unicidade, bloco \lstinline|# (A,B)| \\
Persistência (grupo 10) & bloco \lstinline|# Frame axioms...| \\
Ordem parcial entre metas (grupo 11) & bloco \lstinline|# (ORD)|, com as variáveis \lstinline|ach| e \lstinline|hold|; opção \texttt{--ordem} lida por \lstinline|ler_ordem(...)| \\
Metas da Situação 1 ($S_{f1}$ a $S_{f4}$) & dicionário \lstinline|METAS_SITUACAO1| \\
Plano mínimo (menor $T$ satisfatível) & \texttt{buscar\_plano.py}, laço sobre $T=1,2,\dots$ \\
Verificação independente de P1 a P6 & \texttt{interpretar.py}, função \texttt{validar\_plano} \\
Elos causais e ordem parcial do plano & \texttt{ordem\_parcial.py} \\
$\On(b,y,t)$ derivado & \texttt{interpretar.py}, função \texttt{derivar\_on} \\
\bottomrule
\end{longtable}

\subsection{Diferenças em relação ao manual do professor}
\begin{itemize}
\item O manual exige $\Clr(y)$ na pré-condição de $\Move$. Aqui a exigência foi trocada por $\Free$ por slot (P4), como no item 1. Sem isso, as transições $S_4\to S_5$ das Situações 2 e 3 seriam impossíveis.
\item Foi acrescentada a variável $\Cov$, para que o apoio da estabilidade seja contado no nível certo.
\item O resultado do SAT solver é lido com o arquivo \texttt{.map}, que associa cada número de variável ao seu símbolo, como $\mathit{mv}(d,c,0,0)$.
\end{itemize}

% RASCUNHO da gerencia (Guilherme) para o item 5 -- revisar com Yuri e Gabriel.
\section{Execução Passo a Passo}\label{sec:resultados}

Requisitos: Python 3 e o \texttt{minisat} no PATH. Os arquivos ficam todos na mesma pasta.

\subsection{Uma situação, passo a passo}
\begin{enumerate}
\item \textbf{Gerar o CNF e o mapa.} O primeiro argumento é a situação e o segundo é o horizonte $T$.
\begin{lstlisting}
python3 gerar_cnf.py 3 6
\end{lstlisting}
Saída: \texttt{Situação 3 (T=6): 1267 variáveis, 51308 cláusulas}. Arquivos gerados: \texttt{trab01\_blocos2SAT\_situacao3.cnf} e \texttt{trab01\_blocos2SAT\_situacao3.map}.
\item \textbf{Executar o minisat.}
\begin{lstlisting}
minisat trab01_blocos2SAT_situacao3.cnf resultado3.txt
\end{lstlisting}
O minisat escreve em \texttt{resultado3.txt} a palavra \texttt{SAT} seguida da lista de literais, ou \texttt{UNSAT}.
\item \textbf{Interpretar o resultado.} O script usa o \texttt{.map} correspondente.
\begin{lstlisting}
python3 interpretar.py resultado3.txt --verbose
\end{lstlisting}
\end{enumerate}

\subsection{Plano mínimo automático}
O comando \texttt{python3 buscar\_plano.py 3} repete os passos acima com $T=1,2,3,\dots$ e para no primeiro \texttt{SATISFIABLE}. Esse $T$ é o comprimento mínimo do plano, porque os valores menores deram \texttt{UNSATISFIABLE}. O script grava o \texttt{resultado3.txt} do último $T$ e chama o \texttt{interpretar.py}.

\subsection{Horizontes testados}
O resultado esperado, que é uma propriedade do problema e não do solver, é:
\begin{center}
\begin{tabular}{l l l l}
\toprule
Cenário & \texttt{UNSAT} para & \texttt{SAT} em & Plano mínimo \\
\midrule
Situação 1 ($S_{f3}$) & $T=1$ & $T=2$ & 2 ações \\
Situação 1 ($S_{f4}$) & $T=1,2,3$ & $T=4$ & 4 ações \\
Situação 2 ($S_5$) & $T=1,\dots,4$ & $T=5$ & 5 ações \\
Situação 3 ($S_7$) & $T=1,\dots,5$ & $T=6$ & 6 ações \\
\bottomrule
\end{tabular}
\end{center}
Para $S_{f1}$ e $S_{f2}$ da Situação 1, a busca em largura sobre os 2\,032 estados alcançáveis dá planos mínimos de 8 e 9 ações. Elas são geradas com \texttt{python3 buscar\_plano.py 1 f1} e \texttt{python3 buscar\_plano.py 1 f2}.

\subsection{Com ordem parcial entre metas}
A ordem entre metas entra com a opção \texttt{--ordem}. Na Situação 3, exigir $b(1,1)$ antes de $a(0,1)$ dá $T=5$ e $T=6$ \texttt{UNSAT} e $T=7$ \texttt{SAT}:
\begin{lstlisting}
python3 buscar_plano.py 3 --tmax 8 --ordem='b(1,1)<a(0,1)'
\end{lstlisting}
Os arquivos desse cenário têm o sufixo \texttt{\_ordem}. Já exigir $a(0,1)$ antes de $b(1,1)$ mantém o plano de 6 ações.

\subsection{Arquivos entregues}
Para cada cenário: \texttt{trab01\_blocos2SAT\_situacaoN.cnf}, \texttt{trab01\_blocos2SAT\_situacaoN.map} e \texttt{resultadoN.txt}, com $N=1,2,3$. Os \texttt{resultadoN.txt} são a saída real do minisat para o menor $T$ satisfatível.

% RASCUNHO da gerencia (Guilherme) para o item 5 -- revisar com Yuri e Gabriel.
\section{Interpretação da Saída do SAT Solver}

\subsection{O minisat não conhece nomes simbólicos}
O minisat é um resolvedor puramente booleano. Ele devolve \texttt{SAT} (ou \texttt{UNSAT}) e uma lista de inteiros, onde um número positivo é uma variável verdadeira e um negativo é uma variável falsa. Toda a tradução de número para símbolo vem do arquivo \texttt{.map}, que o \texttt{gerar\_cnf.py} escreve junto com o CNF. Sem o \texttt{.map}, a saída é só uma lista de inteiros sem significado para o domínio.

\subsection{Da saída numérica ao plano em português}
O script \texttt{interpretar.py} executa estes passos:
\begin{enumerate}
\item \textbf{Leitura do mapa:} associa cada número ao seu símbolo, por exemplo \texttt{480} $\to$ \lstinline|move(d,c,0,0)|.
\item \textbf{Filtragem:} mantém os literais positivos. As variáveis \lstinline|move| formam o plano, e as variáveis \lstinline|at| e \lstinline|lev| dão o estado em cada instante.
\item \textbf{Ordenação temporal:} agrupa as ações por $t$ e traduz cada uma para uma frase em português.
\item \textbf{Derivação de $\On$:} no estado final, $\On(b,T)$ se $b$ está no nível 0. Caso contrário, $\On(b,y)$ para todo bloco $y$ no nível $l-1$ cujo $\Span$ se sobrepõe ao de $b$. Isso lista, por exemplo, $d$ sobre $a$ e $b$ em ponte, ou $a$ sobre $c$.
\item \textbf{Verificação independente:} confere cada ação contra as pré-condições P1 a P6 do item 1 e confere se o estado $t+1$ é exatamente o efeito da ação. Essa verificação não usa o CNF, então detecta erro da codificação.
\end{enumerate}

\subsection{Exemplo: Situação 3}
Saída do \texttt{interpretar.py} para um plano mínimo da Situação 3 (6 ações). O minisat pode devolver outro plano com o mesmo número de ações; o que deve coincidir é o tamanho e a validade.
\begin{lstlisting}
PLANO ENCONTRADO (6 ações):
1. t=0: mover bloco 'd' para CIMA de 'c' em p=0
2. t=1: mover bloco 'a' para CIMA de 'b' em p=5
3. t=2: mover bloco 'd' para a MESA em p=2
4. t=3: mover bloco 'a' para CIMA de 'c' em p=0
5. t=4: mover bloco 'b' para CIMA de 'c' em p=1
6. t=5: mover bloco 'd' para a MESA em p=3

ESTADO FINAL (t=6):
  a: ponto inicial p=0, nivel l=1
  b: ponto inicial p=1, nivel l=1
  c: ponto inicial p=0, nivel l=0
  d: ponto inicial p=3, nivel l=0

RELACOES 'on' em t=6:
  a esta sobre: c
  b esta sobre: c
  c esta sobre: MESA
  d esta sobre: MESA

VERIFICAÇÃO INDEPENDENTE (pré-condições P1..P6 do Item 1): plano válido
\end{lstlisting}

\subsection{Comparação com os planos manuais}
A comparação é pelo \emph{tamanho} e pela \emph{validade}, e não pela sequência idêntica, porque um mesmo problema pode ter vários planos mínimos. Por exemplo, na Situação 2 o plano manual leva $b$ para a mesa na primeira ação, e um plano igualmente válido de 5 ações começa com $a$ sobre $d$.
\begin{center}
\begin{tabular}{l c c c}
\toprule
Cenário & Ações no plano manual & $T$ mínimo no SAT & Mesmo tamanho \\
\midrule
Situação 1 ($S_{f3}$) & 2 & 2 & sim \\
Situação 1 ($S_{f4}$) & 4 & 4 & sim \\
Situação 2 ($S_5$) & 5 & 5 & sim \\
Situação 3 ($S_7$) & 6 & 6 & sim \\
\bottomrule
\end{tabular}
\end{center}
Se o SAT desse \texttt{UNSAT} para um $T$ menor do que o plano manual, o resultado seria o esperado, porque o plano manual é um limite superior do menor $T$. Se desse \texttt{SAT} com $T$ menor e o \texttt{interpretar.py} confirmasse o plano como válido, o plano manual não seria mínimo. Se desse \texttt{SAT} e o plano fosse inválido, haveria erro na codificação. Nos quatro cenários da tabela, o menor $T$ do SAT coincide com o do plano manual e o plano devolvido passa na verificação independente.

\paragraph{Exemplo do manual do professor.} O plano de exemplo da seção 6 do manual ($d$ sobre $c$; $a$ para a mesa em $p=4$; $d$ para a mesa em $p=2$; $a$ sobre $c$) não é válido: no terceiro passo, $d$ em $p=2$ cobre os slots 2, 3 e 4, e o slot 4 já está ocupado por $a$ no nível 0. Por isso ele não foi usado como gabarito.

\subsection{Com ordem parcial entre metas}
Com a opção \texttt{--ordem}, o \texttt{interpretar.py} também confere se a ordem foi respeitada: para cada $\varphi_1\prec_P\varphi_2$, em todo instante em que $\varphi_2$ vale, $\varphi_1$ já valeu. O \texttt{ordem\_parcial.py} imprime, para o plano encontrado, os elos causais $A_i\xrightarrow{\varphi}A_j$ e as restrições de ordem necessárias, como na seção sobre ordem parcial.


\end{document}
